\documentclass[11pt]{article}
\usepackage[margin=1in]{geometry}
\usepackage{amsmath,amssymb}
\usepackage{enumitem}
\usepackage[T1]{fontenc}
\usepackage[utf8]{inputenc}
\usepackage{lmodern}
\setlength{\parindent}{0pt}
\setlength{\parskip}{0.5em}
\begin{document}

\section*{IB Physics HL: E1 Problem Collection}
Paraphrased structure-of-the-atom questions from the available 2025 exam pool.
\begin{itemize}[leftmargin=1.25em]
\item
\textbf{Source.} May 2025 HL Paper 1A TZ2, Question 31\\
\textit{Original format: multiple choice.}\\
Using the isotope symbols ${}^{12}_{6}\mathrm{C}$, ${}^{1}_{1}\mathrm{H}$, ${}^{16}_{8}\mathrm{O}$, and ${}^{14}_{7}\mathrm{N}$, consider pure samples of $\mathrm{CO}_2$, $\mathrm{N_2O}$, $\mathrm{CH_4}$, and $\mathrm{H_2O}$. Each sample has the same volume, pressure, and temperature. Determine which sample has the smallest mass.
\item
\textbf{Source.} May 2025 HL Paper 1A TZ2, Question 32\\
\textit{Original format: multiple choice.}\\
Several downward atomic transitions between levels $n=4,3,2,1$ are shown. Determine which transition emits the photon with the longest wavelength.
\item
\textbf{Source.} May 2025 HL Paper 1A TZ2, Question 33\\
\textit{Original format: multiple choice.}\\
In the Geiger-Marsden-Rutherford experiment, alpha particles are accelerated through a potential difference of $7.7\,\text{MV}$ and directed at a thin gold foil. Gold has proton number $79$. Determine the distance of closest approach.
\item
\textbf{Source.} May 2025 HL Paper 1A TZ3, Question 32\\
\textit{Original format: multiple choice.}\\
An energy-level diagram is given with levels $E_1$ to $E_5$, where $E_1$ is the ground state. An electron is initially excited to $E_4$. Determine how many distinct wavelengths can be emitted as electrons cascade back to the ground state.
\item
\textbf{Source.} May 2025 HL Paper 1A TZ3, Question 33\\
\textit{Original format: multiple choice.}\\
Element $X$ has nuclear radius $R$ and nuclear density $\rho$. Element $Y$ has nucleon number eight times that of $X$. Determine the radius and density of the nucleus of $Y$.
\item
\textbf{Source.} November 2025 HL Paper 1A TZ3, Question 33\\
\textit{Original format: multiple choice.}\\
Three statements are made about emission and absorption spectra:
\begin{enumerate}[label=\Roman*., leftmargin=1.5em]
\item They reveal chemical composition.
\item They provide evidence for mass-energy equivalence.
\item They arise from electrons changing energy levels.
\end{enumerate}
Determine which statements are correct.
\item
\textbf{Source.} November 2025 HL Paper 1A TZ3, Question 34\\
\textit{Original format: multiple choice.}\\
For each of several given atoms, imagine removing one electron. Determine for which resulting one-electron system the Bohr model is applicable.
\end{itemize}

\end{document}
